\chapter{Principaux risques identifiés}

\begin{table}[ht]
\centering
\renewcommand{\arraystretch}{1.25}
\small
\begin{tabularx}{\textwidth}{>{\bfseries}c X c c c X}
\toprule
\textbf{ID} & \textbf{Risque} & \textbf{P} & \textbf{I} & \textbf{P×I} & \textbf{Mesure préventive} \\
\midrule
R1 & Manque de temps pour développer toutes les fonctionnalités & 4 & 5 & 20 & Prioriser le MVP et respecter le découpage par phases. \\
R2 & Mauvaise compréhension des besoins réels des étudiants & 3 & 4 & 12 & Interroger quelques étudiants avant de finaliser les exigences et le périmètre. \\
R3 & Difficultés techniques lors de l'intégration & 3 & 4 & 12 & Intégrer régulièrement et bien définir les interfaces entre backend et frontend. \\
R4 & Indisponibilité d'un membre de l'équipe & 3 & 3 & 9 & Partager les connaissances et documenter le travail. \\
R5 & Retard dans la conception de la base de données & 3 & 4 & 12 & Valider tôt le modèle de données et les relations principales. \\
R6 & Annonces fausses ou contenu inapproprié & 3 & 2 & 6 & Mettre en œuvre le signalement, la modération, la suppression d'annonce et le bannissement d'utilisateur. \\
R7 & Dérive du périmètre (scope creep) & 3 & 4 & 12 & Se tenir strictement dans le périmètre du MVP. \\
R8 & Indécision sur le choix technologique (framework / base de données) & 3 & 4 & 12 & Trancher sur les choix technologiques dès la phase de conception. \\
\bottomrule
\end{tabularx}
\caption{Principaux risques identifiés}
\end{table}

\section{Matrice des risques}

\textit{Matrice des risques à compléter ultérieurement.}
